% Figure: the logistic (saturation) growth curve fitted to data, with
% the exponential it agrees with at small population drawn as the
% diverging dashed curve. Carrying capacity K reads off as the upper
% asymptote; the inflection sits at K/2; the early slope matches the
% pure exponential.
\begin{figure}[htbp]
\centering
\resizebox{0.86\linewidth}{!}{%
\begin{tikzpicture}[scale=1.0,
  ax/.style={->, draw=engblack, thick},
  tick/.style={draw=engblack!60, thin},
  curve/.style={draw=engblue, very thick},
  expc/.style={draw=engred, thick, dash pattern=on 4pt off 2pt},
  asy/.style={draw=engblack!55, thick, dashed},
  pt/.style={circle, fill=engblue!75!black, inner sep=1.0pt},
  ann/.style={font=\footnotesize, align=center},
]
  % Horizontal t in [0, 12]; 0.6 cm per unit. Vertical y in [0, 1.2 K],
  % K = 4, scale 0.9 cm per unit.
  \draw[ax] (0, 0) -- (8.0, 0) node[below right, ann] {$t$};
  \draw[ax] (0, 0) -- (0, 5.2) node[left, ann] {population $N$};
  \foreach \x in {2, 4, 6, 8, 10, 12} {
    \draw[tick] ({0.6*\x}, -0.06) -- ({0.6*\x}, 0.06);
    \node[ann, below, yshift=-2pt] at ({0.6*\x}, 0) {\x};
  }
  % Carrying-capacity asymptote K = 4 (y = 0.9*4 = 3.6).
  \draw[asy] (0, {0.9*4}) -- (7.8, {0.9*4});
  \node[ann, anchor=south east, fill=white, inner sep=1pt]
    at (7.8, {0.9*4 + 0.02}) {carrying capacity $K$};
  % Half-capacity / inflection marker K/2.
  \draw[engblack!45, thin, dotted] (0, {0.9*2}) -- (7.8, {0.9*2});
  \node[ann, anchor=south west, fill=white, inner sep=1pt]
    at (0.1, {0.9*2 + 0.02}) {$K/2$ (inflection)};
  % Logistic curve N = K / (1 + ((K-N0)/N0) e^{-r t}), K=4, N0=0.2, r=0.9.
  % => N = 4 / (1 + 19 e^{-0.9 t}). Plot x=t in display units 0.6*t.
  \draw[curve, domain=0:12.8, samples=120, smooth]
    plot ({0.6*\x}, {0.9*( 4/(1 + 19*exp(-0.9*\x)) )});
  % Pure exponential N0 e^{r t} (the small-t agreement, diverging).
  \draw[expc, domain=0:6.0, samples=80, smooth]
    plot ({0.6*\x}, {0.9*( 0.2*exp(0.9*\x) )});
  \node[ann, anchor=south west, text=engred] at ({0.6*4.6}, {0.9*4.4})
    {pure exponential\\(early-time match)};
  % Data points scattered along the logistic.
  \foreach \x in {0.5, 1.5, 2.5, 3.3, 4.0, 4.8, 5.6, 6.6, 8.0, 10.0, 12.0} {
    \pgfmathsetmacro{\yy}{0.9*( 4/(1 + 19*exp(-0.9*\x)) )}
    \node[pt] at ({0.6*\x}, {\yy}) {};
  }
\end{tikzpicture}%
}%
\caption[The logistic curve and the exponential it generalises]{The
logistic growth model $N(t)=K/\bigl(1+((K-N_{0})/N_{0})e^{-rt}\bigr)$
(blue) fitted to a population that grows then saturates. The upper
asymptote reads off the carrying capacity $K$; the inflection sits at
$K/2$, where growth is fastest; the early-time behaviour matches the
pure exponential $N_{0}e^{rt}$ (red dashed) before the saturation term
bends the curve over. The exponential and the logistic agree while
$N\ll K$ and diverge without bound once $N$ approaches $K$, which is
why an exponential fit to the early data over-predicts the late count
by the same mechanism the failure section names for linear
extrapolation of an exponential.}
\label{fig:vol02:ch02:logistic-fit}
\end{figure}
